\section{Warum die Setzungen zuerst kommen}

\textbf{[SETZUNG]} Dieses Dokument steht am Anfang des Sachteils, weil eine
Theorie an der Zahl und Art ihrer Setzungen gemessen wird, nicht an der Zahl
ihrer Ergebnisse.

Die FFGFT behauptet, die Konstanten des Standardmodells aus einem einzigen
dimensionslosen Parameter zu gewinnen. Diese Behauptung ist nur so viel wert,
wie die Liste dessen, was daneben noch vorausgesetzt wird, kurz und offen ist.
Wird die Liste verschwiegen, ist die Ein-Parameter-Aussage wertlos, weil sich
in jede unausgesprochene Voraussetzung beliebig viel Information verstecken
lässt.

Es gibt drei Setzungen. Sie stehen hier vollständig, bevor irgendein Ergebnis
genannt wird.

\section{Erste Setzung: die Zeit-Masse-Dualität}

\textbf{[SETZUNG]} Für die intrinsische Zeitskala $\tilde T$ und die Masse $m$
eines Objekts gilt
\[
  \tilde T \cdot m = 1 .
\]

Das ist keine Definition der Masse und keine Definition der Zeit. Es ist die
Aussage, dass beide Größen zwei Lesarten derselben Sache sind: eine kleine
Masse ist eine lange intrinsische Zeit, eine große Masse eine kurze. Die
Gleichung ist reziprok, also symmetrisch -- keine der beiden Seiten ist die
erklärende.

Nachrichtentechnisch ist das die Fourier-Dualität zwischen Zeit und Frequenz,
in physikalische Sprache übersetzt. Genau von dort stammt sie auch:
die Dualität war der Ausgangspunkt der Theorie, nicht ihr Ergebnis.

\textbf{[B]} Aus dieser einen Gleichung folgt unmittelbar, dass Zeit keine
zusätzliche unabhängige Größe ist. Ein Objekt mit Masse trägt seine Zeitskala
mit sich; es braucht keine von außen angelegte Uhr. Das ist der Grund, warum
die Geometrie mit vier und nicht mit fünf Kreisen auskommt (A030), und der
Grund, warum die Zeit im Zustandsraum liegt und nicht daneben (A060).

\textbf{Was hier gesetzt und nicht gezeigt wird:} dass die Reziprozität exakt
gilt und nicht nur näherungsweise; und dass die Konstante auf der rechten Seite
gleich Eins ist. Die Eins ist eine Einheitenwahl -- jeder andere Wert ließe sich
in die Skalen absorbieren --, aber die Exaktheit der Reziprozität ist eine
inhaltliche Behauptung, die nicht bewiesen, sondern gesetzt wird.

Diese Setzung ist auch der Grund, warum die Theorie sich selbst nicht als Norm
über andere Rahmen stellen kann: $\tilde T\cdot m = 1$ ist eine Festlegung,
keine gefundene Gesetzmäßigkeit. Wer sie nicht teilt, macht keinen Fehler,
sondern arbeitet in einem anderen Rahmen. Die Unterscheidung von Gesetz, Norm
und Setzung wird in A200 ausgeführt; sie beginnt hier.

\section{Zweite Setzung: die kompakte Geometrie}

\textbf{[SETZUNG]} Der zugrundeliegende Raum ist ein vierdimensionaler
kompakter Torus mit einer $Z_3$-Identifikation, geschrieben $T^4/Z_3$.

Drei Teilentscheidungen stecken darin, die getrennt zu bewerten sind.

\textbf{Kompakt statt unendlich.} Das ist die eigentliche Setzung. Sie hat eine
unmittelbare Konsequenz: ein kompakter Raum hat ein diskretes Spektrum. Daraus
folgt, dass Massen als Eigenwerte auftreten und nicht als freie Parameter
eingesetzt werden müssen. Diese Konsequenz ist der ganze Grund für die
Setzung -- und sie ist zugleich der Punkt, an dem die Setzung scheitern kann.

\textbf{Torus statt Kugel.} \textbf{[B]} Diese Wahl ist nicht frei: der Torus
hat eine nichttriviale Fundamentalgruppe, $\pi_1(T^4) \cong \mathbb{Z}^4$, die
Kugel nicht. Windungszahlen -- die diskreten, ganzzahligen, koordinatenfreien
Größen, auf denen der gesamte Informationsbegriff der Theorie ruht (A180) --
existieren auf der Kugel nicht. Wer den Torus durch eine einfach
zusammenhängende Fläche ersetzt, verliert nicht ein Detail, sondern das
tragende Element.

\textbf{Vier Dimensionen.} Hier ist die Buchhaltung unbequem, und das soll sie
bleiben. Dass es \emph{vier} sind und nicht fünf, folgt aus der ersten Setzung:
weil $\tilde T\cdot m = 1$ die Zeit an die Masse bindet, braucht es keinen
eigenen fünften Kreis für die Zeit. Insofern ist die Vier abgeleitet. Dass die
erlebte Struktur aber gerade $3+1$ ist -- drei Raumrichtungen und eine Zeit --
ist \textbf{empirische Eingabe} und wird von der Theorie nicht hergeleitet.
Ältere Formulierungen im Korpus waren an dieser Stelle zu stark; die Korrektur
ist im Register als P34 verzeichnet und hier eingearbeitet.

\textbf{Die $Z_3$-Identifikation.} \textbf{[SETZUNG]} Die Dreizähligkeit wird
gesetzt. Sie liefert die dreielementige Faser, auf der die drei Generationen
als Spektrum eines zyklischen Operators erscheinen (A110). \textbf{[B]} Was
sich zeigen lässt: die Dreizähligkeit ist kristallographisch zulässig, die
Fünfzähligkeit nicht -- die Bedingung $2\cos(2\pi/n) \in \mathbb{Z}$ erlaubt
$n \in \{1,2,3,4,6\}$. Das begründet nicht, warum es drei sein müssen; es
schließt aber eine ganze Klasse von Alternativen aus.

\textbf{[SETZUNG]} Die Richtung ist dabei festzuhalten: \emph{festgelegt} wird
$n=3$ nicht durch die Symmetrie, sondern durch die beobachtete Zahl der
Generationen -- die Beobachtung wählt die Symmetrie, nicht umgekehrt. $Z_3$ wird
gesetzt, motiviert durch die Sparsamkeit der dreipunktigen Konfiguration und im
Einklang mit der Generationenzahl. Die Umkehrung -- \enquote{drei Generationen
sind keine freien Parameter, sondern die drei erlaubten $Z_3$-Fixpunkte} -- liest
die Beobachtung als Folge der Symmetrie und ist zu vermeiden.
Eine Setzung wird durch das Vorliegen von Gründen nicht zur Ableitung.

\section{Dritte Setzung: der Zahlenwert}

\textbf{[SETZUNG]} Der Parameter ist
\[
  \xi = \frac{4}{30000} = \frac{1}{3\cdot 2^2\cdot 5^4}
  \approx 1{,}333\cdot 10^{-4}.
\]

Diese Setzung ist die angreifbarste der drei, und sie zerfällt in zwei Teile
mit ganz verschiedenem Status. Die Trennung ist wichtig genug, um sie
ausdrücklich zu machen -- ältere Darstellungen im Korpus haben beide Teile
gemeinsam als hergeleitet geführt, was nicht haltbar war (Register P33).

\textbf{Die Form ist strukturell.} \textbf{[B]} Die Faktoren $4/3$, die $3$ und
die $2^2$ haben eine geometrische Herkunft: sie stammen aus der
Dimensionsstruktur und der Dreizähligkeit. Insoweit ist die Gestalt des
Ausdrucks nicht gewählt.

\textbf{Die Größenordnung ist es nicht.} Der Faktor $5^4$ -- also die Frage,
warum $\xi$ bei $10^{-4}$ liegt und nicht bei $10^{-2}$ oder $10^{-6}$ --
ist \emph{nicht} vorwärts hergeleitet. Er wird auf zwei Weisen begründet, und
beide sind ehrlich als das zu bezeichnen, was sie sind: intuitiv über
harmonische Verhältnisstrukturen, und empirisch über den Higgs-Sektor. Keine
der beiden Begründungen ist eine Ableitung aus der Geometrie.

\textbf{[SETZUNG]} Damit gilt: der Zahlenwert von $\xi$ ist eine Eingabe. Die
Theorie behauptet nicht, ihn zu erzeugen; sie behauptet, dass aus ihm alles
Weitere folgt.

Eine Warnung gehört an dieselbe Stelle, weil sie den häufigsten Denkfehler
betrifft. Jede dimensionslose Größe lässt sich als Potenz von $\xi$ schreiben,
wenn man den Exponenten frei wählen darf. Ein Ausdruck der Form
$X = \xi^N$ ist für sich \emph{aussagelos}. Er wird erst zu einer Aussage,
wenn der Exponent unabhängig festgelegt ist -- durch die Geometrie, durch eine
Abzählung, durch eine Rekursionstiefe. Wo das in dieser Serie nicht möglich
ist, wird es gesagt.

\section{Was aus den drei Setzungen folgt}

\textbf{[B]} Die drei Setzungen zusammen legen fest: einen kompakten Raum mit
diskretem Spektrum, eine dreielementige Faser, eine Skala. Mehr nicht. Alles
Weitere in dieser Serie ist entweder aus diesen dreien abgeleitet, oder es ist
ausdrücklich als weitere Setzung markiert.

Die Probe darauf ist die Zählung selbst. Wenn im Verlauf der Serie eine vierte
oder fünfte Setzung nötig wird, steht sie an ihrem Ort mit der Marke
\textbf{[SETZUNG]} -- und sie muss dann in die Bilanz von A230 eingehen. Eine
Theorie mit drei offen gelegten Setzungen ist besser als eine mit einer
verschwiegenen; eine mit sieben offen gelegten wäre schlechter als eine mit
drei. Die Buchführung entscheidet, nicht die Rhetorik.

\section{Der Toleranzmaßstab an dieser Stelle}

In diesem Dokument steht bewusst keine einzige Vergleichszahl. Der Grund ist
methodisch: Setzungen haben keine Messgenauigkeit. Es ist sinnlos zu fragen, ob
$\tilde T\cdot m = 1$ auf sechs Stellen stimmt -- die Gleichung wird gesetzt,
nicht gemessen.

Prüfbar sind Setzungen nur mittelbar, über ihre Konsequenzen, und dort greift
der Maßstab aus A010: eine Konsequenz gilt als getragen, wenn ihr Rest
innerhalb der Bestimmungstoleranz bleibt und nicht durch einen freien Parameter
geschlossen werden muss. Die ersten solchen Zahlen stehen in A100.

\section{Prüfskripte}

Zwei Aussagen dieses Dokuments sind rechnerisch und daher hinterlegt:

\begin{itemize}
  \item die kristallographische Restriktion $2\cos(2\pi/n)\in\mathbb{Z}$ mit
    dem Ergebnis $n\in\{1,2,3,4,6\}$ und dem Ausschluss der Fünfzähligkeit:
    \texttt{python/A\_Serie\_Skripte/a020\_kristallographie.py}
  \item die Aussagelosigkeit von $X=\xi^N$ bei freiem Exponenten -- gezeigt an
    beliebig gewählten Zielgrößen, für die sich stets ein passendes $N$ finden
    lässt: \texttt{python/A\_Serie\_Skripte/a020\_xi\_potenz\_trivialitaet.py}
\end{itemize}

\section{Was offen bleibt}

\textbf{Erstens ist die Zahl drei nicht bewiesen minimal.} Dass es genau diese
drei Setzungen braucht, ist eine Behauptung über die Sparsamkeit der
Konstruktion. Es ist nicht gezeigt, dass sich keine davon aus den anderen
gewinnen lässt, und ebenso wenig, dass zwei genügen würden.

\textbf{Zweitens ist die Dreizähligkeit nicht begründet, nur zugelassen.} Die
kristallographische Restriktion schließt die Fünfzähligkeit aus und erlaubt
$\{1,2,3,4,6\}$. Warum von den erlaubten Werten gerade die Drei realisiert ist,
sagt die Geometrie nicht.

\textbf{Drittens ist die Größenordnung von $\xi$ nicht nur intuitiv begründet,
aber die Ableitung ist noch nicht in der A-Serie geführt.} Der Exponent $-4$
folgt aus $D_f=3-\xi$ und der Vier-Dimensionalität des Universums; $\kappa=7$
emergiert als einzige konsistente Lösung des e-p-$\mu$-Systems; der Faktor
$4/3$ ist durch den Casimir-Effekt unabhängig bestätigt. Diese Herleitung steht
in A015. Was \emph{noch} fehlt: eine Ableitung der Größenordnung aus den drei
Setzungen selbst, ohne Rückgriff auf gemessene Massen als Eingang. Bis dahin
trägt die quantitative Seite einen gemessenen Anker.

\textbf{Viertens ist $3+1$ eine Eingabe.} Sie ist als solche markiert und wird
in keinem späteren Dokument der Serie stillschweigend zur Ableitung
aufgewertet. Ob das durchgehalten wird, prüft das Audit-Skript.

\section{Ersetzt}

Dieses Dokument nimmt auf: Dok.~003 (Grundlagen), Dok.~008
($\xi$ und $e$), Dok.~009 ($\xi$-Ursprung), Dok.~041
(Parameterherleitung), Dok.~042 ($\xi$-Parameter Teilchen), Dok.~253
($\xi$-Zahlenraum). Eingearbeitete Registereinträge: P33 (Status der
Magnitude), P34 (Dimensionsstruktur), P35 ($\xi^N$-Trivialität), R56 ($Z_3$ als
Setzung; Richtung Beobachtung$\to$Symmetrie, nicht umgekehrt).

\section{Verwendung von KI-Werkzeugen}

Dieses Dokument ist unter Verwendung KI-gestützter Werkzeuge entstanden. Die
Fragestellungen, die theoretischen Setzungen und alle inhaltlichen
Entscheidungen stammen vom Autor; die Werkzeuge formulieren aus, rechnen nach,
prüfen auf Widerspruch und erstellen Prüfskripte. Sie treffen keine
inhaltlichen Entscheidungen. Die Verantwortung für Inhalt und Fehler liegt
vollständig beim Autor. Der ausführliche Wortlaut steht in A010.
